% =========================================================
\chapter{Récapitulatif général pour l'examen}
% =========================================================

\begin{rembox}[title=Comment utiliser ce chapitre]
Ce chapitre rassemble \textbf{tout le cours} en quelques pages : le fil conducteur, le
formulaire complet, les méthodes pas à pas et les pièges classiques. Il ne contient rien de
nouveau : chaque ligne renvoie à un chapitre (Cxx). La veille de l'examen, lire d'abord la
boîte \og chèvre\fg{}, puis le formulaire, puis les méthodes.
\end{rembox}

\begin{goatbox}[title=Tout le cours expliqué à une chèvre]
L'électricité, c'est de l'eau : la \textbf{tension} est la hauteur, le \textbf{courant} le
débit (C01). Une \textbf{résistance} est un tuyau étroit, un \textbf{condensateur} un
réservoir, une \textbf{inductance} une lourde roue à eau (C02). L'eau ne disparaît jamais aux
carrefours, et faire le tour de la montagne ramène à la même altitude : ce sont les deux lois
de \textbf{Kirchhoff} (C02). Vue de l'extérieur, toute machine compliquée faite de ces pièces
ressemble à une seule pile avec un seul tuyau (\textbf{Thévenin}) ou une seule pompe avec une
fuite (\textbf{Norton}), et on peut étudier chaque source séparément puis additionner
(\textbf{superposition}, C03). Quand l'eau va et vient comme une vague (\textbf{sinus}), les
réservoirs et les roues deviennent des \og tuyaux\fg{} dont l'étroitesse dépend de la vitesse
de la vague : ce sont les \textbf{impédances}, et le rapport sortie/entrée est la
\textbf{fonction de transfert} (C04). On la dessine avec des règles droites sur du papier
logarithmique : \textbf{Bode} (C05). Quand on ouvre le robinet d'un coup (\textbf{créneau}),
le réservoir se remplit en exponentielle, 63\,\% après $\tau$, presque tout après $5\tau$
(C06). Enfin la \textbf{diode} est un portillon à sens unique qui demande 0,7\,V pour
s'ouvrir, et la \textbf{Zener} cède aussi dans l'autre sens à $V_Z$ (C07).
\end{goatbox}

\section{Le fil conducteur du cours}

\begin{center}
\renewcommand{\arraystretch}{1.3}
\begin{tabular}{p{0.08\linewidth}p{0.3\linewidth}p{0.52\linewidth}}
\textbf{Ch.} & \textbf{Thème} & \textbf{Ce qu'il apporte} \\ \hline
C00 & Introduction & mesurer/décider/agir ; analogique vs numérique ; linéarisation \\
C01 & Signaux et sources & $U$, $I$, conventions, sources réelles, moyenne, efficace \\
C02 & R, L, C, Kirchhoff & lois des composants, mailles/n\oe uds, Millman, Kennelly, équivalences \\
C03 & Méthodes d'analyse & dipôles, Thévenin, Norton, superposition \\
C04 & Impédances & complexes, $\underline Z$, fonction de transfert, filtres \\
C05 & Bode & forme canonique, dB, asymptotes \\
C06 & Impulsions & saut indiciel, continuité de $V_C$, formule magique \\
C07 & Diodes & non-linéaire, droite de charge, modèles, Zener \\
\end{tabular}
\end{center}

\section{Formulaire complet}

